% v25. Exact factorization of the original degenerate pair residual.
\section{退化面の積表示と互いに素な領域の排除}
\label{dg25:section}

三高係数が同時に非負である canonical 整数模型の記号
\eqref{hd23:data} を用いる．退化面 $a=uh-B(F+1)=0$ では，
元の pair 残余を単なる一次式として解く以上に，正整数の積として扱える．
これは $x\ge1$ の残余に対する新しい互いに素な場合の障害と，
Schur 整数 $b$ だけを固定した有限復元を与える．
退化面全体の不存在は結論しない．

\subsection{元の残余の正整数積への変換}

\begin{theorem}[退化面の積表示]
\label{dg25:product}
$a=0$ とし，
\begin{equation}\label{dg25:data}
 \begin{gathered}
 A=B(s-1)-1=Bs-b,\quad R=b^2-sT,\qquad
 L=(F+1)R-hA,\quad M=uR-BA
 \end{gathered}
\end{equation}
と置く．このとき $L,M$ は正整数であり，
\begin{equation}\label{dg25:identities}
 hM=BL,\qquad FL=bF+s(T-h),\qquad
 \mathfrak r=M(t-y)-tT.
\end{equation}
従って元の pair 方程式 $\mathfrak r=0$ は
\begin{equation}\label{dg25:positive-product}
 Mj=tT,\qquad j=t-y\in\{1,\ldots,t\}
\end{equation}
を要求する．特に $y<t$，$M\ge T$，$M\mid tT$ である．
\end{theorem}

\begin{proof}
既知の退化関係より $s\ell=bF$ であり，
$K_0=h(s-1)+1$ と置いたとき
\[
 R(F^2+F+1)=(bK_0-sh)F+b^2-sh.
\]
ここで $bK_0-sh=hA+b$ である．従って
\begin{equation}\label{dg25:Q}
 Q:=R(F^2+F+1)-(hA+b)F-b^2+sh=0.
\end{equation}
定義から $hM-BL=Ra=0$ であり，$R=b^2-sT$ を用いると
$FL-bF-s(T-h)=Q=0$ となる．

$M$ の正値性は恒等式として確認できる．$m=F^2+F+1$，
\[
 C=F\{b(t-1)+s\}+b^2t+bs-b>0
\]
と置けば，$t=u-s+1$ を使った直接展開により
\begin{equation}\label{dg25:positive-M}
 mM-C=uQ+(FA-s)a
\end{equation}
である．ゆえに $M=C/m>0$ であり，$hM=BL$ と $h,B>0$ から $L>0$ でもある．

退化面上では $T$ は $y$ に依存せず，
\[
 S=FT-B(t-y)-\ell,\quad
 \alpha=F+tz+uy,\quad z=\ell-u
\]
である．これらを元の残余に代入すると，より一般の恒等式
\begin{equation}\label{dg25:original-identity}
 -(s+tb)S+(s\alpha-t)T-\alpha b
 =M(t-y)-tT+(s\ell-bF)(tT+1)
\end{equation}
を得る．$s\ell=bF$ により主張の残余式になる．
模型ではその残余が零であり $M,t,T>0$ なので $j>0$ である．
$y\ge0$ より $j\le t$，従って $M=tT/j\ge T$ となる．
\end{proof}

差方程式だけから式~\eqref{dg25:positive-product} を結論してはいけない．
最後に $\mathfrak r=0$ を用いることが不可欠である．
一般の三高模型に $j>0$ を仮定するのも正しくない．ここでは退化面上で証明した．

\begin{corollary}[互いに素な領域の無上界排除]
\label{dg25:coprime}
同じ退化模型で $h\ge b$ なら，
\begin{equation}\label{dg25:small-L}
 1\le L\le b-1,\qquad
 g:=\gcd(h,b-1)\ge2,\qquad
 \frac hg\le b-1,\qquad
 T\le M\le b-2
\end{equation}
である．従って $h\ge b$ かつ $\gcd(h,b-1)=1$ の領域は存在しない．
特に $x\ge1$ の全ての退化模型にこの結論が適用される．
\end{corollary}

\begin{proof}
退化面の既知の範囲は $R\ge1$ である．従って
$T=(b^2-R)/s<b^2/s\le h$ となる．最後は $s\ge h\ge b$ による．
式~\eqref{dg25:identities} から $L<b$ であり，正整数性と合わせて最初の範囲を得る．
$h\mid BL$ は $h/g\mid L$ と同値だから，$h/g\le b-1$ である．
$h\ge b$ の下では $g=1$ は不可能である．
さらに $M=BL/h\le B^2/h<B$ なので，整数性から $M\le b-2$ を得る．

$x\ge1$ の場合には，退化関係 $s-bx>0$ から $s\ge b+1$ と
$x\le(s-1)/b$ が従う．
\[
 s-b-\frac{s-1}{b}=\frac{(b-1)(s-b-1)}b\ge0
\]
なので $x\le s-b$，従って $h=s-x\ge b$ である．
\end{proof}

\subsection{正の第一係数を持つ退化面の固定 \texorpdfstring{$b$}{b} 復元}

\begin{corollary}[固定 $b$ に対する完全有限手順]
\label{dg25:fixed-b}
$a=0,x\ge1$ の三高非負係数模型は $b\ge3$ を満たす．
固定した $b$ ごとに，次の有限手順で全て復元できる．
\begin{enumerate}
 \item $s=b+1,\ldots,b^2-1$，$x=1,\ldots,\lfloor(s-1)/b\rfloor$ を取り，
 $h=s-x$，$B=b-1$，$A=Bs-b$，$g=\gcd(h,B)$ とする．
 \item $T=1,\ldots,\min\{b-2,\lfloor(b^2-1)/s\rfloor\}$ と，
 $h/g$ の正倍数 $L\le B$ を取る．$R=b^2-sT$，$M=BL/h$ とし，$M\ge T$ を検査する．
 \item
 \[
 F=\frac{s(h-T)}{b-L},\quad u=\frac{B(F+1)}h,\quad
 \ell=\frac{bF}s,\quad t=u-s+1,\quad z=\ell-u
 \]
 が整数で，$t\ge1,z\ge0$，$F=B+ux+sz$，
 $L=(F+1)R-hA$ を満たすものを残す．
 \item $M\mid tT$ を検査し，$y=t-tT/M\ge0$ とする．
 元の尺度復元~\eqref{hd23:reconstruction} で $S,T,U$ を計算し，
 三尺度の正整数性，$d\ge1$，$\alpha\ge k\ge b+1$，六つの符号条件，
 全ての行和・pair products と指定 $D$ を確認する．
 \item 通過した模型に完全 matching の有理正 weights を与える．
\end{enumerate}
この手順では $b$ 自体に上界を証明していない．$x=0$ の全退化面には
系~\ref{ax24:degenerate-finite} の固定 $s$ 手順を使う．
\end{corollary}

\begin{proof}
$b\ge3$ は $1\le T\le b-2$ から従う．$s\ge b+1$ は前の証明による．
$R\ge1,T\ge1$ より $s\le b^2-1$ である．
残りの全範囲と復元式は前定理と前系から従い，$b-L>0$ なので零分母はない．
従って全ての所望の模型が列挙される．
最後に元の全方程式と正値性を検査するため，逆向きの十分性も明白である．
これは必要条件だけを満たす整数を模型と認定する手順ではない．
\end{proof}

\paragraph{検算の射程．}
\nolinkurl{verify_degenerate_divisor_v25.py} は式~\eqref{dg25:positive-M}，
\eqref{dg25:original-identity} と二つの積恒等式を整数多項式として係数ごとに検査する．
\texttt{--fixed-b b} は上の証明された有限範囲を全て処理し，任意の分母や $t,z,y$ の
探索上限を追加しない．有限個の $b$ に対する実行結果を，退化面全体の不存在へ外挿しない．
